% part: intuitionistic-logic
% chapter: propositions-as-types
% section: normalization

\documentclass[../../../include/open-logic-section]{subfiles}

\begin{document}

\olfileid{int}{pty}{nor}

\olsection{सामान्यीकरण}

या विभागात आपण सिद्ध करू की विशिष्ट न्यूनीकरण-क्रम वापरल्यास
कोणतीही निष्पत्ती सामान्य निष्पत्तीत न्यूनीत करता येते. यालाच
\emph{सामान्यीकरणाचा गुणधर्म} म्हणतात. यासाठी आपण विधाने-हे-प्रकार
अनुरूपता वापरू: प्रत्येक सिद्धता-पद सामान्य रूपात न्यूनीत करता
येते हे दाखवू; त्यावरून नैसर्गिक निगमनाच्या
!!{derivation}s चे सामान्यीकरण मिळेल.

प्रथम पदांची गुंतागुंत मोजणारी काही फलने परिभाषित करू.
!!a{formula}s ची \emph{लांबी}~$\len{!A}$ पुढीलप्रमाणे:
\begin{align*}
  \len{p} &= 0\\
  \len{!A \land !B} &= \len{!A} + \len{!B} + 1\\
  \len{!A \lor !B} &= \len{!A} + \len{!B} + 1 \\
  \len{!A \lif !B} &= \len{!A} + \len{!B} + 1.
\end{align*}
रेडेक्स~$M$ ची गुंतागुंत त्याच्या \emph{कट-कोटी}ने,
म्हणजे~$\cutrank{M}$ ने मोजतो:
\begin{align*}
  \cutrank{(\lambd[\typeof{x}{!A}][\typeof{N}{!B}])Q} &=
  \len{!A} + \len{!B} + 1 \\
  \cutrank{\ande{i}{\andi{\typeof{M}{!A}}{\typeof{N}{!B}}}} &=
  \len{!A} + \len{!B} + 1 \\
  \cutrank{\ore{\ori{i}{!A_{i}}{\typeof{M}{!A_i}}}
    {\typeof{x_1}{!A_1}}{\typeof{N_1}{!C}}
    {\typeof{x_2}{!A_2}}{\typeof{N_2}{!C}}} &=
  \len{!A_1} + \len{!A_2} + 1
\end{align*}
सिद्धता-पदाची गुंतागुंत त्यातील सर्वाधिक गुंतागुंतीच्या
रेडेक्सने मोजतो; पद सामान्य असेल, तर ती~$0$ असते:
\[
  \maxrank{M} = \max \{\cutrank{N} | N \text{ हे उपपद आहे } M \text{ आणि रेडेक्स आहे}\}
\]

\begin{lem}\ollabel{lem:subst}
  $\Subst{M}{\typeof{N}{!A}}{\typeof{x}{!A}}$ हा रेडेक्स असेल
  आणि $M \not\ident x$ असेल, तर पुढीलपैकी एक शक्यता असते:
  \begin{enumerate}
  \item $M$ स्वतः रेडेक्स आहे; किंवा
  \item $M$ चे रूप $\ande{i}{x}$ आणि $N$ चे रूप
    $\andi{P_1}{P_2}$ आहे;
  \item $M$ चे रूप $\ore{i}{x_1}{P_1}{x_2}{P_2}$ आणि $N$ चे
    रूप $\ori{i}{}{Q}$ आहे;
  \item $M$ चे रूप $x Q$ आणि $N$ चे रूप
    $\lambd[x][P]$ आहे.
  \end{enumerate}
  पहिल्या बाबतीत $\cutrank{\Subst{M}{N}{x}} = \cutrank{M}$;
  इतर बाबतींत $\cutrank{\Subst{M}{N}{x}} = \len{!A}$.
\end{lem}
\begin{proof}
  ~$M$ वर आगमनाने सिद्ध करू.
  \begin{enumerate}
  \item $M$ हे एकच चर $y$ असेल आणि $y \not\ident x$ असेल,
    तर $\Subst{y}{N}{x}$ हे $y$ आहे; म्हणून ते रेडेक्स नाही.
  \item $M$ चे रूप $\andi{N_1}{N_2}$, किंवा $\lambd[x][N]$,
    किंवा $\ori{i}{!A}{N}$ असेल, तर
    $\Subst{M}{\typeof{N}{!A}}{\typeof{x}{!A}}$ चेही तेच
    रूप असते; म्हणून ते रेडेक्स नसते.
  \item $M$ चे रूप $\ande{i}{P}$ असेल, तर दोन बाबी पाहू.
    \begin{enumerate}
      \item $P$ चे रूप $\andi{P_1}{P_2}$ असेल, तर
        $M \ident \ande{i}{\andi{P_1}{P_2}}$ हा रेडेक्स आहे;
        आणि स्पष्टच,
        \[
        \Subst{M}{N}{x} \ident
        \ande{i}{\andi{\Subst{P_1}{N}{x}}{\Subst{P_2}{N}{x}}}
        \]
        हाही रेडेक्स आहे. दोघांच्या कट-कोटी समान आहेत.
      \item $P$ एकच चर असेल, तर प्रतिस्थापनाने रेडेक्स
        मिळण्यासाठी तो~$x$ असलाच पाहिजे, आणि~$N$ चे रूप
        $\andi{P_1}{P_2}$ असले पाहिजे. आता
        \[
        \Subst{M}{N}{x} \ident \Subst{\ande{i}{x}}{\andi{P_1}{P_2}}{x},
        \]
        विचारात घ्या. हे $\ande{i}{\andi{P_1}{P_2}}$ आहे.
        त्याची कट-कोटी $\cutrank{\ande{i}{\andi{P_1}{P_2}}}$
        म्हणजेच $\len{!A}$ आहे.
    \end{enumerate}
  \end{enumerate}
  $\ore{N}{x_1}{N_1}{x_2}{N_2}$ आणि $P Q$ या बाबीही
  अशाच प्रकारच्या आहेत.
\end{proof}

\begin{lem}
  $M$ चे $M'$ मध्ये संकुचन होत असेल आणि~$M$ च्या प्रत्येक
  योग्य रेडेक्स उपपद~$N$ साठी
  $\cutrank{M} > \cutrank{N}$ असेल, तर
  $\cutrank{M} > \maxrank{M'}$.
\end{lem}

\begin{proof}
  बाबींनुसार सिद्ध करू.
  \begin{enumerate}
  \item $M$ चे रूप $\ande{i}{\andi{M_1}{M_2}}$ असेल, तर
    $M'$ हे $M_i$ आहे. $M_i$ चे कोणतेही उपपद हे~$M$ चे
    योग्य उपपदही असल्यामुळे दावा सिद्ध होतो.
  \item $M$ चे रूप
    $(\lambd[\typeof{x}{!A}][N])\typeof{Q}{!A}$ असेल,
    तर $M'$ हे $\Subst{N}{\typeof{Q}{!A}}{\typeof{x}{!A}}$
    आहे. $M'$ मधील रेडेक्स घ्या. एकतर~$N$ मध्ये त्याला
    अनुरूप व समान कट-कोटीचा रेडेक्स आहे; गृहीतकामुळे ती
    कट-कोटी $\cutrank{M}$ पेक्षा कमी आहे. अन्यथा त्याची
    कट-कोटी $\len{!A}$ आहे; व्याख्येनुसार ती
    $\cutrank{(\lambd[x^{!A}][N])Q}$ पेक्षा कमी आहे.
  \item $M$ चे रूप
    \[
    \ore{\ori{i}{}{\typeof{N}{!A_i}}}
        {\typeof{x_1}{!A_1}}{\typeof{N_1}{!C}}
        {\typeof{x_2}{!A_2}}{\typeof{N_2}{!C}},
    \]
    असेल, तर $M' \ident \Subst{N_i}{N}{\typeof{x_i}{!A_i}}$.
    ~$M'$ मधील रेडेक्स घ्या. एकतर~$N_i$ मध्ये त्याला
    अनुरूप व समान कट-कोटीचा रेडेक्स आहे; गृहीतकामुळे ती
    कट-कोटी $\cutrank{M}$ पेक्षा कमी आहे. अन्यथा त्याची
    कट-कोटी $\len{!A_i}$ आहे; व्याख्येनुसार ती
    $\cutrank{\ore{\ori{i}{}{\typeof{N}{!A_i}}}
      {\typeof{x_1}{!A_1}}{\typeof{N_1}{!C}}
      {\typeof{x_2}{!A_2}}{\typeof{N_2}{!C}}}$ पेक्षा कमी आहे.
  \end{enumerate}
\end{proof}

\begin{thm}
  सर्व सिद्धता-पदे सामान्य रूपात न्यूनीत होतात;
  सर्व !!{derivation}s सामान्य !!{derivation}s मध्ये न्यूनीत होतात.
\end{thm}

\begin{proof}
  दुसरा दावा पहिल्यावरून मिळतो. पहिला दावा सिद्धता-पद~$M$
  साठी $m=\maxrank{M}$ वर पूर्ण आगमनाने सिद्ध करू.
  \begin{enumerate}

  \item $m=0$ असेल, तर $M$ आधीच सामान्य आहे.
  \item
    अन्यथा,~$M$ मधील कट-कोटी~$m$ असलेल्या रेडेक्सांची
    संख्या~$n$ यावर आगमन करू.
    \begin{enumerate}
    \item $n=1$ असेल, तर प्रत्येक योग्य रेडेक्स उपपद~$P$
      यासाठी $m = \cutrank{N} > \cutrank{P}$ असणारा
      कोणताही रेडेक्स~$N$ निवडा. असे रेडेक्स अस्तित्वात
      असलेच पाहिजे, कारण कोणत्याही पदाची उपपदे सांत असतात.

      $N$ च्या संकुचनफलाला~$N'$ म्हणा. आता प्रमेयिकेने
      $\maxrank{N'} < \maxrank{N}$; म्हणून
      $\cutrank=m$ असलेल्या रेडेक्सांची संख्या~$n$ कमी होते.
      त्यामुळे~$m$ कमी होतो ($1$ किंवा अधिकने), आणि~$m$
      वरील आगमन-गृहीतक वापरता येते.
    \item आगमनाच्या पायरीत $n>1$ माना. प्रक्रिया तशीच
      आहे; फरक इतकाच की~$n$ कमी होऊनही धन राहतो, त्यामुळे
      $m$ बदलत नाही. म्हणून~$n$ वरील आगमन-गृहीतक वापरतो.
    \end{enumerate}
\end{enumerate}
\end{proof}

प्रकारयुक्त $\lambd$-पदांचे सामान्यीकरण होते याचा अर्थ
अशा प्रत्येक पदाला \emph{सामान्य रूप मिळते}. $\lambd$-पदांचे
सामान्य रूपापर्यंत $\beta$-न्यूनीकरण करणे म्हणजे संगणन चालवणे,
आणि सामान्य रूप म्हणजे त्याचे मूल्य, असे मानल्यास प्रकारयुक्त
$\lambd$-कलनातील प्रत्येक संगणन अखेरीस थांबून मूल्य देते.
शिवाय प्रकार-संरक्षणामुळे त्या मूल्याचा प्रकार योग्य असतो.
उदाहरणार्थ,~$N$ चा प्रकार $!A \lif !B$ असेल
(म्हणजे ते प्रकार~$!A$ च्या आर्ग्युमेंटवर चालणारे आणि
प्रकार~$!B$ चे मूल्य देणारे फलन असेल) आणि प्रकार~$!A$
च्या~$M$ या पदावर (निवेशावर) ते लावले, तर $(NM)$ चे
$N'$ या पदात (निर्गमात) न्यूनीकरण होते आणि त्या निर्गमाचा
प्रकार~$!B$ असेल याची हमी मिळते.

प्रत्यक्षात प्रकारयुक्त $\lambd$-कलनातील पदे उपपदांवर
$\beta$-न्यूनीकरण लावण्याच्या फक्त एका पद्धतीने सामान्य रूपात
जात नाहीत; $\beta$-न्यूनीकरण लावण्याची \emph{कोणतीही}
पद्धत अखेरीस सामान्य रूपात थांबते. या गुणधर्माला
\emph{प्रबळ सामान्यीकरण} म्हणतात. याचा अर्थ केवळ
संगणन थांबून मूल्य मिळेल याची खात्री करणारी एक पद्धत
आहे एवढाच नाही; संगणन चालवण्याची \emph{कोणतीही}
पद्धत अखेरीस मूल्य देते.

संगणन चालवण्याच्या वेगवेगळ्या पद्धतींनी \emph{तेच}
मूल्य मिळते हे कसे कळते? आतापर्यंत (प्रकार-संरक्षणामुळे)
एवढेच माहीत आहे की मिळालेल्या सर्व मूल्यांचा प्रकार एकच
असतो. प्रकारयुक्त $\lambd$-कलनाला आणखी एक महत्त्वाचा
गुणधर्म आहे: ते \emph{संगमशील}, म्हणजे त्याला
\emph{चर्च–रॉसर गुणधर्म} आहे. $N \red N_1$ आणि
$N \red N_2$ असतील, तर एखादे पद~$N'$ असे असते की
$N_1 \red N'$ आणि $N_2 \red N'$. लक्षात ठेवा की
$N \red N'$ याचा अर्थ $\redone$-न्यूनीकरणाच्या
शून्य किंवा अधिक पायऱ्यांनी~$N'$ मिळते.
$N_1$ सामान्य रूपात असेल, तर $N_1 \red N'$ असणारे
एकमेव पद~$N'$ म्हणजे~$N_1$ स्वतःच (शून्य
$\redone$-न्यूनीकरण पायऱ्यांनी). त्यामुळे~$N_1$ आणि
$N_2$ दोन्ही सामान्य रूपात असतील, तर $N_1 = N' = N_2$.
दुसऱ्या शब्दांत, $\redone$ प्रबळ सामान्यीकारी असल्यामुळे
पद न्यूनीत करण्याची प्रत्येक पद्धत अखेरीस सामान्य रूप देते.
$\red$ ला चर्च–रॉसर गुणधर्म असल्यामुळे कोणतीही दोन
सामान्य रूपे समान असतात. म्हणून प्रकारयुक्त
$\lambd$-कलनातील संगणन नेहमी थांबते आणि ते कसेही
चालवले तरी मिळणारे मूल्य (सामान्य रूप) तेच असते.

संबंध संगमशील आहे हे सिद्ध करणे सहसा अवघड असते.
तथापि, पुढील दुर्बल अट सहज सिद्ध करता येते:

\begin{prop}[दुर्बल चर्च–रॉसर गुणधर्म]
$N \redone N_1$ आणि $N \redone N_2$ असतील, तर एखादे
पद~$N'$ असे असते की $N_1 \red N'$ आणि $N_2 \red N'$.
\end{prop}

\begin{proof}
  $N$ मधील रेडेक्स~$M_1$ किंवा~$M_2$ त्याच्या संकुचनफलाने
  बदलल्यास अनुक्रमे~$N_1$ किंवा~$N_2$ मिळतात.
  $M_1$ आणि~$M_2$ हे~$N$ चे उपपद असल्याने त्यांचे
  अंशतः छेदन होऊ शकत नाही. म्हणजे एक रेडेक्स दुसऱ्याचे
  उपपद असेल, किंवा ते~$N$ मधील वेगवेगळ्या न छेदणाऱ्या
  स्थानी असतील. दुसरी बाब घ्या:~$N$ चे रूप $N(M_1,M_2)$
  आहे. $M_1$ च्या जागी त्याचे संकुचनफल~$M_1'$ ठेवल्यास
  $N(M_1',M_2)$ हे~$N_1$ मिळते; $M_2$ च्या जागी त्याचे
  संकुचनफल~$M_2'$ ठेवल्यास $N(M_1,M_2')$ हे~$N_2$
  मिळते. दोघांचेही $N(M_1',M_2')$ मध्ये न्यूनीकरण होते.
  दुसऱ्या बाबतीत~$M_2$ हे~$M_1$ चे उपपद आहे असे माना
  (उलट बाब सममित आहे), आणि~$M_1$ कोणत्या प्रकारचा
  रेडेक्स आहे यानुसार बाबी वेगळ्या करा. उदाहरणार्थ,
  $M_1 \ident \proj{1}{\pair{O_1}{O_2}}$ असेल, तर त्याचे
  $O_1$ मध्ये न्यूनीकरण होते. $M_2$ हे~$O_1$ चे उपपद
  असेल, तर~$M_2$ चे परिवर्तन केल्यावर~$O_1'$ मिळते.
  म्हणून आधी~$M_1$ आणि नंतर~$M_2$ चे परिवर्तन केल्यावर
  $O_1'$ मिळते. पण आधी~$M_2$ चे परिवर्तन केल्यावर
  $\proj{1}{\pair{O_1'}{O_2}}$ मिळते, आणि त्याचेही
  $O_1'$ मध्ये न्यूनीकरण होते.
\end{proof}

\begin{lem}[न्यूमनची प्रमेयिका]
$\redone$ प्रबळ सामान्यीकारी असेल आणि त्याला दुर्बल
चर्च–रॉसर गुणधर्म असेल, तर तो संगमशील आहे.
\end{lem}

\begin{proof}
  $\redone$ प्रबळ सामान्यीकारी असल्यामुळे प्रत्येक पूर्ण
  न्यूनीकरण-क्रम सामान्य रूपात संपतो. सामान्य रूपे एकमेव असणे आणि
  $\red$ संगमशील असणे या सममूल्य अटी आहेत. म्हणून एखाद्या
  पद~$N$ ला~$N'$ आणि~$N''$ ही दोन वेगळी सामान्य रूपे
  आहेत असे माना; ती पुढील न्यूनीकरण-क्रमांनी मिळतात:
  \begin{align*}
  N & \redone N_1' \redone N_2' \redone \dots \redone N_k' = N' \text{ आणि}\\
  N & \redone N_1'' \redone N_2'' \redone \dots \redone N_l'' = N''
  \end{align*}
  (न्यूनीकरण-क्रमाची लांबी $>0$ असलीच पाहिजे, अन्यथा
  $N$ आधीच सामान्य रूपात असेल.)

  $N_1' = N_1''$ असेल, तर त्याला दोन वेगळी सामान्य रूपे
  ($N'$ आणि $N''$) आहेत. $N_1' \neq N_1''$ असेल, तर
  दुर्बल चर्च–रॉसर गुणधर्मामुळे एखादे~$N'''$ असे आहे की
  $N_1' \red N'''$ आणि $N_1'' \red N'''$.
  $N' \neq N''$ असल्यामुळे $N''' \neq N'$ किंवा
  $N''' \neq N''$ यांपैकी किमान एक अट खरी आहे.
  निर्णायक मुद्दा असा: प्रबळ सामान्यीकरणामुळे या सामायिक
  पदाला एक सामान्य रूप मिळते. ते मूळ दोन वेगळ्या सामान्य
  रूपांपैकी किमान एकापेक्षा वेगळे असते. त्यामुळे~$N_1'$
  किंवा~$N_1''$ यांपैकी एकाला दोन वेगळी
  सामान्य रूपे आहेत. आपण दाखवले की~$N$ ला दोन वेगळी
  सामान्य रूपे असतील, तर एखादे~$N_1$ असे आहे की
  $N \redone N_1$ आणि~$N_1$ लाही दोन वेगळी सामान्य रूपे
  आहेत. ही प्रक्रिया पुढे चालू ठेवून
  $N \redone N_1 \redone N_2 \redone \dots$ असा क्रम
  मिळेल; पण $\redone$ प्रबळ सामान्यीकारी असल्यामुळे
  हे अशक्य आहे.
\end{proof}

\end{document}
